\section{Effective positive endpoint and temporal attachment}
\label{sec:effective-positive}
The positive endpoint must supply constants before an overlap time can
be chosen. We now make the full-saddle field, mixed collars and actual
mean bounds effective. The order is field comparison, local geometry,
regularized primitives, then the complete tuple. This is the content of
Note~54, Sections~A--B \cite{n54}; the compact extension required after
choosing the overlap is treated in Section~\ref{sec:all-time}.

\subsection{A finite field theorem on the full complex buffer}
Retain the implicit saddle and determinant in
\eqref{eq:full-saddle-model}, and write
\[
 \mathcal L=\log(2+t+|x|),\quad \mathcal N=1+|x|+t\mathcal L,
 \quad e=\frac{t^{3/2}}{\mathcal N^2},\quad d=\frac{t}{\mathcal N^3}.
\]
\[
 A_\Sigma=|S_t((1+iz)/2)|+|S_t((1-iz)/2)|.
\]
\begin{theorem}[Effective all-height positive field]
\label{thm:effective-positive-field}
For $t\ge T_F:=\exp(4096)$, the original saddle branch is well defined
on $|\Im z|\le6$, at every real height. On this strip,
\[
 |H-P|\le10000e A_\Sigma,\qquad
 |P_t+P_{zz}|\le8192d A_\Sigma.
\]
On $|\Im z|\le5$,
\begin{equation}
 |\partial_z^j(H-P)|\le C_j e\mathcal L^j A_\Sigma,
 \qquad C_j=2^{16}j!8^j,\quad 0\le j\le3.
 \label{eq:ep-field}
\end{equation}
All constants are independent of the subsequently chosen core and source.
\end{theorem}
\begin{proof}
Put $w=\sigma+i\tau_0$, $-3\le\sigma\le4$, and first take
$\tau_0\ge0$. The gamma remainder used in
Section~\ref{sec:effective-inputs} gives, for $\Re p\ge2$ and
$m=|p|\ge4096$,
\[
 \alpha(p)=\tfrac12\Log(p/(2\pi))+r(p),\quad |r|\le2/m,
 \quad |\alpha^{(j)}|\le j!/m^j\ (1\le j\le4),
 \quad |r'|\le24/m^2.
\]
For all $m$ in this half-plane the weaker $|r|\le4/m$ also holds.
On the rectangle
\[
 \sigma+t\mathcal L/8\le\Re p\le\sigma+t\mathcal L,
 \qquad \tau_0\le\Im p\le\tau_0+t,
\]
the map $p\mapsto w+t\alpha(p)/2$ has Lipschitz constant at most
$20/\mathcal L<1/2$, real increment in
$[t\mathcal L/5,t\mathcal L/3]$, and imaginary increment in $[0,t/2]$.
For the lower imaginary sign, if $a=\Re p>48$, then
$\Re\alpha'(a+i\tau)\ge(a/2-24)/(a^2+\tau^2)>0$; integrate
vertically from the real axis. Thus the fixed point is in the strict
interior, and
\begin{equation}
 \mathcal N/10\le|p|\le2\mathcal N,\quad
 a_*:=\Re(p-w)\in[t\mathcal L/5,t\mathcal L/3],\quad
 |\Im(p-w)|\le t/2,\quad |B-1|\le20/\mathcal L<1/8.
 \label{eq:ep-branch}
\end{equation}
Conjugation handles the other sign. Nonvanishing of $B$ supplies analytic
neighborhoods; uniqueness identifies overlapping charts and the earlier
exterior branch.

Use the exact reflected-arm decomposition of the heat convolution in
Section~\ref{sec:full-saddle}. Each arm starts at $v_0=2-\sigma$
($|v_0|\le5$) and is moved horizontally through the same saddle, retaining
its oriented vertical connector. The length-three middle interval is
also retained. Close each finite contour before taking its right end to
infinity; quadratic Gaussian decay dominates the gamma growth there.
For the shifted variable $u\ge-a_*/2$, the Taylor integral remainder
and its quadratic part obey
\[
 |\Gamma_p(u)|\le\frac{3u^2}{2m},\qquad
 |E_p(u)|\le\frac{3|u|^3}{m^2}.
\]
The modulus of both exponential integrands is at most
$e^{-u^2/(2t)}$. Interpolating their exponents and using
$\int_\R |u|^3e^{-u^2/(2t)}du=4t^2$ bounds the central error by
$24t^{3/2}/m^2\le2400e$ relative to its arm carrier.

For the left piece the all-size gamma bound gives
\[
 \int_{A-h}^{A}\Re\alpha(v+i\tau)\,dv
 \ge h\Re\alpha(p)-D(A),\qquad
 D(A)=A/2+4\log(A/2)+4\le3A,
\]
where $A=\Re p$. The connector uses $|\Im\alpha(2+i\tau)|<2$,
and the middle uses
$\sup_{-1\le c\le2}|\xi(c+i\tau)/8|
 \le225(1+|\tau|)|M(2+i\tau)|$.
The resulting complete ledger, per arm where appropriate, is
\begin{center}\small
\begin{tabular}{@{}ll@{}}\toprule
Piece & Relative upper allowance\\\midrule
Central quadratic replacement & $2400e$\\
Missing Gaussian and actual infinite ends & $e^{-t\mathcal L^2/200}$ each\\
Left horizontal piece, connector, middle & $e^{-t\mathcal L^2/400}$ each\\
All non-unit arithmetic modes & $128e^{-t\mathcal L/10+t/2}
                                      \le e^{-t\mathcal L/20}$\\\bottomrule
\end{tabular}\end{center}
Indeed $-a_*^2/(4t)+D(A)\le-t\mathcal L^2/200$; the connector
adds at most $t+1$, and the largest prefactor is $e^{2\mathcal L}$.
At the declared onset these losses fit the weaker displayed exponent.
Also $\mathcal N\le e^{\mathcal L}\mathcal L$ and
$e\ge e^{-3\mathcal L}$, so every exponential remainder is at most $e$.
Two arms and the middle fit the constant $10000$. No finite arithmetic
truncation or omitted infinite end is used.

For $D=(\log S_t)_w$, exact differentiation gives
$|D|\le2\mathcal L$, $|D_w|\le3/m$, $|\Im D|\le2$.
On a disk of radius $(8\mathcal L)^{-1}$ the carrier and $e$ each
change by at most a factor two. Cauchy's formula therefore gives
\eqref{eq:ep-field}; strip six pays complete disks centered in strip
five. The exact determinant residual formula in
Section~\ref{sec:positive-evolution}, with \eqref{eq:ep-branch}, gives
$8192d$. In particular no asymptotic error bound is differentiated.
\end{proof}

\subsection{Acquire mixed collars before dividing by a field}
Write $P=2\mathcal A\cos\Theta$ with compatible analytic logarithms.
The same exact formulas yield
\begin{equation}
 |(\log\mathcal A)_z|\le2,\quad
 \mathcal L/8\le\Theta_x\le\mathcal L\quad(y=0),\quad
 |\Theta_{zz}|\le10/\mathcal N.
 \label{eq:ep-phase}
\end{equation}
In particular $\Re\Theta_z>0$ throughout the full buffer. Integrating
vertically shows that $\Im\Theta$ has the sign of $y$; the model zeros
there are real and simple.

Let the exact-heat source field $G$ satisfy
\begin{equation}
 |\partial_z^j(G-H)|\le e\mathcal L^j A_\Sigma/16,
 \qquad j\le3,
 \label{eq:ep-source-premise}
\end{equation}
on the used core and its complete strip-four disks. For every
$F_s=(1-s)G+sP$, $0\le s\le1$, use the disks about a real $P$ root
$c$ with radii $(16\mathcal L_c)^{-1}$ and
$(64\mathcal L_c)^{-1}$. Equation~\eqref{eq:ep-phase} gives
$|P(z)|\ge\mathcal A(c)\mathcal L_c|z-c|/32$ on these disks
and an off-collar floor $2^{-20}A_\Sigma$.
Rouch\'e's theorem on the two boundaries gives one simple real root
$a_s$ and no additional zeros on the consumed support, with
\[
 |a_s-c|\le2^{40}e/\mathcal L_c,\qquad
 2^{-10}\mathcal A(c)\mathcal L_c
 \le |(F_s)'(a_s)|\le2^{10}\mathcal A(c)\mathcal L_c.
\]
Reality follows from conjugation and uniqueness. On the outer disk
factor $F_s=(z-a_s)b_s$. Boundary comparison and the maximum principle
for $b_s$ and $1/b_s$ bound the units; on the half outer disk,
\[
 |\log(b_s/b_H)|\le2^{62}e,\qquad
 |\partial_z\log(b_s/b_H)|\le2^{70}e\mathcal L_c.
\]
The pole residue needs the exact phase formula, rather than the coarse
individual $O(\mathcal L)$ unit derivative:
\[
 v_P=2\mathcal A_z/\mathcal A+\Theta_{zz}/\Theta_z,\quad |v_P|\le8.
\]
Comparison at the displaced roots using jets through three gives
\begin{equation}
 \frac{(F_s)_{zz}}{F_s}=\frac{v_s}{z-a_s}+g_s,\quad
 |v_s|\le16,\quad |v_s-v_H|\le2^{100}e\mathcal L_c,
 \quad |g_s-g_H|\le2^{100}e\mathcal L_c^2.
 \label{eq:ep-poles}
\end{equation}
The regular-part bound follows on a larger boundary after subtracting
both poles, then by the maximum principle inside. All disks used here
lie in the source-controlled strip four. These claims are restricted to
that domain, and impose no earlier central simplicity premise.

\subsection{The temporal primitive and the physical weights}
For the holomorphic core put $\Delta=P-G$, $F=(1-v)G+vP$ and
\[
 T_m[a]=\int_\Omega a\,w_F^m
                 \frac{|\Delta\overline F|}{|F|^2+q},\qquad 0\le m\le2.
\]
Take a full enlarged support of area $A$, at most $n_\omega\omega$
collars of radius $r_0/\omega$, $\omega\asymp\mathcal L$, with
$|F|\ge c_s A_X\omega|z-a_F|$ on collars and
$|F|\ge c_g A_\Sigma,c_g A_X$ off them. Let
$|\Delta|\le D_T e A_X$ on collars, absorbing carrier comparisons in
$D_T$. Direct integration of the simple pole gives
\begin{equation}
 T_m[1]\le D_T\left(A/c_g+
              \frac{2\pi n_\omega r_0}{c_s\omega_0}\right)e.
 \label{eq:ep-temporal-primitive}
\end{equation}
Here $\omega\ge\omega_0>0$. Every intersecting collar is enlarged to a
whole disk before integration. This proves the unweighted $O(e)$ bound
without an inverse regulator.

For $m\ge1$ put $\Lambda=\sup_\Omega\sqrt q/A_X$ and
$a_*=\sup_\Omega a$. On a collar,
$w_F\le h^2/(|z-a_F|^2+h^2)$ with $h=\Lambda/(c_s\omega)$.
The radial integral $\int_0^\infty h^2/(u^2+h^2)du=\pi h/2$ gives
\begin{equation}
 T_m[a]\le a_*D_T e\left(
       A\Lambda^2/c_g^3+\frac{\pi^2 n_\omega\Lambda}{c_s^2\omega}\right).
 \label{eq:ep-weighted-primitive}
\end{equation}
The unweighted estimate is used for $Z_0[1]$; the spatially weighted
terms of \eqref{eq:tc-signed-heat} retain $Z_1,Z_2$.
Replacing those by $Z_0$ would lose the necessary suppression.

The physical source saddle $L_s$ supplies a finite mass anchor.
On $|u-L_s|\le\kappa^{-1/2}$ the source exponent decreases by at
most two, since $\kappa\ge6t+36$ and the explicit third derivative
remainder is bounded there. The first source mode is at least half its
leading term and the others are positive. Hence
\[
 H(t,0)\ge M_0(t)/32,\qquad
 H(t,0)\le3A_0(t),\qquad M_0(t)/A_0(t)\le96,
 \quad A_0=S_t(1/2)>0.
\]
Together with $|(\log\mathcal A)_z|\le2$ this proves, for a fixed
buffer of radius eight and $X/\kappa>128$,
\begin{equation}
 \Lambda\le96\kappa^{-1}
 e^{-(X-8)^2/(2\kappa)+2X}\le96e^{-X}.
 \label{eq:ep-mass-suppression}
\end{equation}
On $X/\kappa\le128$ the exact physical coefficient table
\eqref{eq:positive-exact-regulator-jets} bounds the spatial weights by
fixed constants and the time weights by fixed constants times
$\mathcal L^2$. On its complement each is a stated polynomial in
$1+X$. For any integer $n\ge0$,
\[
 (1+X)^n e^{-cX}\le\sum_{j=0}^n\binom njj!c^{-j}.
\]
This absorbs every weighted exterior term. The bounded ratio used in
this argument is $L_s/\kappa$; the height-dependent
$\mathcal L/\kappa$ is not bounded on all heights.

\subsection{Actual covariance and a noncircular overlap choice}
The unit replacement and pole translation arguments of
Lemma~\ref{lem:regularized-poles} give $s,w\le C e$ and
$v_m\le C e\mathcal L^2$ for arbitrary positive pointwise $q$.
Pole translation retains $e(1+|\log e|)$; here
$1+|\log e|\le5\mathcal L$.
For actual means use
\[
 F_z/F=(z-a_F)^{-1}+b_F'/b_F,
 \qquad |\nabla S_F|\le2|F_z/F|+w_F|\nabla r|.
\]
Whole-collar integration and \eqref{eq:ep-mass-suppression} bound its
integral by $C\mathcal L$. Moving the remaining derivatives to $K$
therefore bounds every actual mean through order three before applying
the covariance product rule. On a spatial band add $F_{\bar z}/F$
using the proved gap floor. The same cap applies to a support meeting
both core and band. A core primitive must never be multiplied by a
compact coefficient depending on the as-yet unchosen upper time.

Here is one explicit common payment. Normalize units by
$A_X\omega$, take $\omega=\mathcal L$, $A=n_\omega=128$,
$r_0=3/128$, and all lower geometric constants $2^{-64}$.
All upper geometric and jet constants can be enlarged to $2^{128}$.
The preceding radial formulas, coefficient weights and factorial sums
then give $C_{\rm raw}=2^{4096}$ for the raw primitives and actual
gradient cap. The largest required exterior weight degree is finite
(at most sixteen after absorption); each factorial sum fits $2^{128}$.
The signed-stress identity is applied before any absolute square bound,
and the actual global mean product is subtracted as in
Section~\ref{sec:compact-transfer}.

Let $a_i$ be the kernel constants in
Section~\ref{sec:effective-inputs}. Define
\begin{equation}
 K_*=1+5\sum_{i+j\le4}a_i a_j,\qquad
 C_+=2^{65536}(1+128+K_*)^{64}.
 \label{eq:ep-Cplus}
\end{equation}
Substitution in the finite primitive maps, including
\eqref{eq:tc-signed-heat}, gives each holomorphic scalar row the bound
$C_+(1+c_b/\Delta_t)e\mathcal L^2$.
The maps use sums, products, and one bilinear mean-product row; eight
raw factors and sixty-four kernel/regime factors already overcount them.
Thus \eqref{eq:ep-Cplus} is an explicit enlargement, independent of
$T_+,\Delta_t,R_c$ and the source. Compact spatial-band costs retain
their separate reserve.

Since $e\mathcal L^2\le t^{-1/2}$, choose
\begin{equation}
 \Delta_t=1,\quad \mu_{\rm coll}=2^{-256},\quad C_{\rm jet}=2^{29},
 \qquad
 T_+=\max\left\{128,e^{4096},
   [8C_+(1+c_b)/\theta]^2,[4C_{\rm jet}/\mu_{\rm coll}]^2\right\}.
 \label{eq:ep-Tplus}
\end{equation}
The onset also pays the root/unit smallness gates, whose multipliers
are at most $2^{256}$. The positive core and pure model tail now cost
at most $\theta/8$ per scalar. For a finite core and $b=T_++1$, an
explicit full-complex carrier floor is
\begin{equation}
 A_{\rm late}=\exp[-4b\log^2(b+2)-2(R_c+5)].
 \label{eq:ep-late-floor}
\end{equation}
Indeed $|\mathcal A(t,z)|\ge A_0(t)e^{-2(R_c+5)}$; the gamma
integral on $[1,2]$ and \eqref{eq:ep-branch} give
$A_0(t)\ge e^{-4b\log^2(b+2)}$. This poor floor is used only to
choose a finite source accuracy. It was never inverted in the temporal
primitive or in $C_+$.
